%==============================================================================
% Exemplo 2
%==============================================================================

%------------------------------------------------------------------------------%
\begin{question}
  Considerando a equação
  \[
    x^2 + 8xy + y^2 - 15 = 0
  \]
  \vspace{-\baselineskip}
  \begin{enumerate}
    \item escreva a equação na forma matricial
    \item diagonalize a matriz da parte quadrática
    \item faça a mudança de variáveis
    \item identifique a cônica
  \end{enumerate}
\end{question}

%------------------------------------------------------------------------------%
\begin{resolution}{Equação Matricial}
  Equação na forma matricial
  \[
    v^tAv+B^tv+f=0
  \]
  \[
    v=\begin{bmatrix}x\\y\end{bmatrix}
    \qquad
    A=\begin{bmatrix}
      1&4\\
      4&1
    \end{bmatrix}
  \]
  \[
    B=\begin{bmatrix}0\\0\end{bmatrix}
    \qquad
    f=-15
  \]
\end{resolution}

%------------------------------------------------------------------------------%
\begin{resolution}{Diagonalizando $A$ -- Polinômio Característico}
  \begin{align*}
    p(\lambda)
    & = \det(A-\lambda I)
    \\
    & = \begin{vmatrix}
      1-\lambda&4\\
      4&1-\lambda
    \end{vmatrix}
    \\
    & = (1-\lambda)^2-16
    \\
    & = \lambda^2-2\lambda-15
  \end{align*}
\end{resolution}

%------------------------------------------------------------------------------%
\begin{resolution}{Diagonalizando $A$ -- Autovalores}
  Equação característica
  \[
    \lambda^2-2\lambda-15=0
  \]
  \[
    (\lambda-5)(\lambda+3)=0
  \]
  Autovalores
  \[
    \lambda_1=5
    \qquad
    \lambda_2=-3
  \]
\end{resolution}

%------------------------------------------------------------------------------%
\begin{resolution}{Autovetores Associados a $\lambda_1=5$}
  \begin{align*}
    (A-5I)v_1 & = \vec{0} \\
    \begin{bmatrix}
      -4&4\\
      4&-4
    \end{bmatrix}
    \vetor{x}{y}
    & = \vetor{0}{0}
  \end{align*}
  Portanto
  \[
    -4x+4y=0
    \qquad
    y=x
  \]
  Escolhendo \(x=1\)
  \[
    v_1=\vetor{1}{1}
  \]
\end{resolution}

%------------------------------------------------------------------------------%
\begin{resolution}{Autovetores Associados a $\lambda_2=-3$}
  \begin{align*}
    (A+3I)v_2 & = \vec{0} \\
    \begin{bmatrix}
      4&4\\
      4&4
    \end{bmatrix}
    \vetor{x}{y}
    & = \vetor{0}{0}
  \end{align*}
  Portanto
  \[
    4x+4y=0
    \qquad
    y=-x
  \]
  Escolhendo \(x=1\)
  \[
    v_2=\vetor{1}{-1}
  \]
\end{resolution}

%------------------------------------------------------------------------------%
\begin{resolution}{Matriz Ortogonal}
  Normalizando os autovetores
  \[
    u_1=\frac{1}{\sqrt{2}}\vetor{1}{1}
    \qquad
    u_2=\frac{1}{\sqrt{2}}\vetor{1}{-1}
  \]
  Matriz ortogonal
  \[
    P=\frac{1}{\sqrt{2}}
    \begin{bmatrix}
      1&1\\
      1&-1
    \end{bmatrix}
  \]
  Matriz diagonal
  \[
    D=P^tAP
    =\begin{bmatrix}
      5&0\\
      0&-3
    \end{bmatrix}
  \]
\end{resolution}

%------------------------------------------------------------------------------%
\begin{resolution}{Mudança de Variáveis}
  Mudança de variáveis
  \[
    v=Pw
    \qquad
    w=\vetor{\mu}{\nu}
  \]
  \[
    \begin{bmatrix}
      x\\y
    \end{bmatrix}
    =
    \frac{1}{\sqrt{2}}
    \begin{bmatrix}
      1&1\\
      1&-1
    \end{bmatrix}
    \begin{bmatrix}
      \mu\\\nu
    \end{bmatrix}
  \]
  Portanto
  \[
    x=\frac{\mu+\nu}{\sqrt{2}}
    \qquad
    y=\frac{\mu-\nu}{\sqrt{2}}
  \]
\end{resolution}

%------------------------------------------------------------------------------%
\begin{resolution}{Equação Transformada}
  Como \(B=\vec{0}\), temos \(E=P^tB=\vec{0}\).

  \begin{align*}
    w^tDw+E^tw+f & = 0
    \\
    5\mu^2-3\nu^2-15 & = 0
    \\
    5\mu^2-3\nu^2 & = 15
  \end{align*}
\end{resolution}

%------------------------------------------------------------------------------%
\begin{resolution}{Identificação}
  Dividindo a equação por \(15\)
  \[
    \frac{\mu^2}{3}-\frac{\nu^2}{5}=1
  \]
  A cônica é uma hipérbole com centro na origem.

  Semieixos
  \[
    a=\sqrt{3}
    \qquad
    b=\sqrt{5}
  \]

  O eixo real está na direção de \(\mu\).
\end{resolution}

%------------------------------------------------------------------------------%
\begin{resolution}{Esboço}
  \centering
  \includegraphics{I-08-Rotacao-ex-3}
\end{resolution}
%------------------------------------------------------------------------------%
